% ===================================================================
%  કોષ્ટક નં. ૮ — લણણી. --- શિકાર, દ્રાક્ષની લણણી, માછીમારી.
%
%  Traduction, faite en 2026, en gujarati d'aujourd'hui, sur l'ido de
%  texto/io. Regles de la colonne : voir 10-kostak-01.tex. Deux
%  scenes, deux numerotations.
%
%  ECART DECLARE, ET IL FAUT LE SAVOIR AVANT D'ECRIRE, NON APRES. Le
%  bloc t08-c2-03-1 est l'un des trois que renvoji.py laisse passer
%  sans rien dire : le fac-simile ido y compose
%  « \textsuperscript{13)} », malforme, sans la parenthese ouvrante,
%  et la machine ne voit donc pas ce renvoi. Un ecart declare rend le
%  bloc AVEUGLE — le controle n'y regarde plus rien du tout, pas
%  meme ce qu'il verifie partout ailleurs. L'ordre juste, pris sur le
%  volume francais, est 3, 12, 14, 13, 15, 16 ; il a ete pose a la
%  main et recompte a la main.
%
%  LA FAUX (25) N'EST JAMAIS VENUE EN INDE, LA FAUCILLE (12) N'EN EST
%  JAMAIS PARTIE. C'est la trouvaille de ce tableau et elle est nette,
%  parce que la planche grave les deux outils cote a cote. « દાતરડું »
%  est la faucille, le mot est vieux, l'objet est celui que toute
%  l'Inde tient encore dans la main droite au moment de la moisson.
%  La faux — la grande lame au bout d'un manche a deux poignees, qu'on
%  balance debout — n'a pas de nom gujarati, parce qu'elle n'a pas
%  d'exemplaire gujarati : on moissonne accroupi ici, pas debout. On
%  ecrit « લાંબા હાથાનું દાતરડું », la faucille a long manche, et le
%  fleau (40) « ઝૂડવાની લાકડી », le baton a battre. C'EST LA LOI DU
%  HAOUSSA VERIFIEE SUR UN AUTRE CONTINENT : ce qui a pousse ici se
%  nomme ici, ce qui n'est jamais venu se decrit.
%
%  ET LE SEIGLE (43) EST LE PIEGE DE CE TABLEAU. Le mot qui tombe sous
%  la main est « રાઈ » — et « રાઈ » est la MOUTARDE, que le Gujarat
%  cultive par trains entiers et dont la fleur jaune couvre le
%  Saurashtra en janvier. Le seigle n'a pas de nom ici parce qu'il n'a
%  pas de champ ici. On ecrit « રાય », l'anglais transcrit, qui ne se
%  confond avec rien. Le haoussa n'avait pas de mot du tout pour le
%  meme numero ; le gujarati en a un de trop, ce qui est pire.
%
%  CINQUIEME REFUS DE LA COLONNE : LE COQUELICOT (2). « ખસખસ » est le
%  pavot, tout le monde le connait, on en met la graine dans les
%  patisseries — mais c'est le pavot somnifere, la plante du champ
%  licite, celle de l'opium. Le coquelicot de la planche est la
%  mauvaise herbe rouge du champ de ble, qui n'a jamais pousse ici et
%  ne porte pas ce nom-la. On ecrit « પોપી ».
%
%  MAIS LE NENUPHAR (18) EST LE CAS INVERSE, DANS LE MEME TABLEAU. Le
%  francais n'a qu'un mot, nenuphar, pour deux plantes que le gujarati
%  separe depuis toujours : « કમળ » est le lotus, celui qui sort de
%  l'eau sur sa tige et qui est la fleur nationale de l'Inde, et
%  « પોયણું » est le nymphea, celui qui flotte a plat. La planche
%  grave des feuilles posees sur l'eau. On ecrit « પોયણાં ».
%
%  ET LE CRAPAUD (1) ET LA GRENOUILLE (17) SONT UN SEUL MOT. Le meme
%  tableau qui vient de couper en deux le nenuphar francais fond en un
%  les deux betes de l'ido : « દેડકો » les nomme toutes les deux, et
%  le gujarati n'a rien pour les separer. La colonne ecrit « દેડકા »
%  au (1) et « નાનો દેડકો » au (17), et c'est la planche qui fait la
%  difference, pas la langue. UNE LANGUE NE DECOUPE PAS LE MONDE
%  PARTOUT PLUS FIN NI PARTOUT PLUS GROS : elle le decoupe fin la ou
%  elle regarde. Ici on regarde l'eau, pas ce qui saute dedans.
%
%  ENFIN LA BERGERONNETTE (19) EST « ધોબણ », LA LAVANDIERE. L'oiseau
%  bat de la queue au bord de l'eau comme la femme bat le linge sur la
%  pierre, et c'est de la que le gujarati le nomme. Le francais dit
%  bergeronnette, la petite bergere ; les deux langues ont vu une
%  femme au travail, et pas la meme.
% ===================================================================

%%K t08-noto-1 noto
\VUnotes{143.2mm}{%
(*) કેમ કે આ શબ્દ ચોક્કસ નથી : શણની લણણી છે, દ્રાક્ષની લણણી, કે સફરજનની
લણણી ? કદાચ « \VUgras{mesono} » વાપરવું સારું રહેશે, જે Mondo XIV, પાન
૨૪૨ માં સૂચવાયું છે. પણ આજ સુધી અકાદમીએ આ નવું મૂળ સ્વીકાર્યું નથી અને તે
રૂઢ ઇડિસ્ટ નિયમ પ્રમાણે ચર્ચાની રાહ જુએ છે.%
}

%%K t08-apar-1 apar
\VUsaut{5.0mm}
\VUtitre{15.0pt}{60}{કોષ્ટક નં. ૮}
\VUsaut{3.0mm}
\VUcentre{13.2pt}{90}{\textit{પ્રથમ દૃશ્ય.}}
\VUsaut{3.0mm}
\VUpk{10.2pt}{40}{લણણી (*).}
\VUsaut{4.0mm}
\VUpk{10.2pt}{40}{ગ્રામ્ય પ્રદેશનું રૂપ.}
\VUsaut{3.0mm}

%%K t08-c1-01-1 p
1. --- \VUgras{ઉનાળા} સાથે ગ્રામ્ય પ્રદેશનું રૂપ ફરી એક વાર બદલાયું.
ચાસમાં સોંપાયેલું બી ઊગ્યું ; લીલો નાનો છોડ જમીનમાંથી બહાર આવ્યો અને તેને
કણસલું બેઠું. દાણો બંધાયો, પછી તે પાક્યો, અને હવે તો માત્ર લણવાનું જ બાકી
છે.

%%K t08-c1-02-1 p
2. --- \VUgras{ખેતરનાં ફૂલ} ખીલી ઊઠ્યાં છે. \VUgras{કોર્નફ્લાવર}
\textsuperscript{(1)}, \VUgras{પોપી} \textsuperscript{(2)} અને
\VUgras{ડિજિટાલિસ} \textsuperscript{(3)} ઘઉંનાં ખેતરોના એકસરખા રંગમાં
પોતાની તાજી \VUgras{પાંખડીઓનો} આનંદી વિરોધ ભેળવે છે.

%%K t08-c1-03-1 p
3. --- ફળઝાડ પાંદડાંથી ભરાઈ ગયાં ; ફળ પાક્યાં. નાનું \VUgras{ચેરીનું ઝાડ}
\textsuperscript{(4)} સુંદર \VUgras{ચેરીથી} \textsuperscript{(5)} લચી
પડ્યું છે, જેને છોકરાં ખૂબ આનંદથી તોડે છે અને ખાય છે.

%%K t08-c1-04-1 p
4. --- હવેથી ઢોર આખો દિવસ ખુલ્લી હવામાં ગાળી શકે છે.

%%K t08-c1-04-2 p
\VUgras{ઘેટાંનાં બચ્ચાં} \textsuperscript{(6)} મોટાં થયાં અને સુંદર
\VUgras{ઘેટીઓ} \textsuperscript{(7)} તથા સુંદર \VUgras{મેંઢા}
\textsuperscript{(8)} બન્યાં, જાડા ઊનના પડવાળાં. તેઓ શાંતિથી
\VUgras{ગોચરનું} \textsuperscript{(9)} \VUgras{ઘાસ} ચરે છે, જે
\VUgras{ડેઝીનાં ફૂલોથી} ટપકાંવાળું છે.

%%K t08-c1-04-3 p
જલદી જ \VUgras{આ} પ્રાણીઓને કાતરવાનું શરૂ થશે, જેથી તેમનું \VUgras{ઊન}
મળે, જેમાંથી આપણાં કપડાંનું કાપડ બને છે.

%%K t08-tit-2 sub
\VUsaut{5.0mm}
\VUpk{10.2pt}{40}{ભરવાડ.}
\VUsaut{3.4mm}

%%K t08-c1-05-1 p
5. --- \VUgras{ભરવાડ} \textsuperscript{(10)} તેમની બહુ સંભાળ લેતો હોય
એવું લાગતું નથી. તેને તો માત્ર ક્ષિતિજ પર થઈને જતા તોફાનની ચિંતા છે, અને
\VUgras{ઘેટાં ચરાવનારો} \textsuperscript{(13)}, જે પોતાની
\VUgras{તુંબડી} \textsuperscript{(14)} હોઠે લગાડે છે, તે તો વળી વધારે
બેફિકર લાગે છે.

%%K t08-c1-06-1 p
6. --- ગરમી અસહ્ય છે ; કેમ કે \VUgras{કૂતરો} \textsuperscript{(15)} પણ
પોતાની તરસ છિપાવવા આવે છે ; તે \VUgras{નાની નદીનું} \textsuperscript{(16)}
ઠંડું પાણી ચાટે છે અને કાંઠાના ઘાસમાં લપાયેલા બિચારા \VUgras{નાના દેડકાને}
\textsuperscript{(17)} ડરાવે છે. તે સ્વચ્છ પાણીમાં કૂદી પડે છે, જ્યાં
\VUgras{પોયણાં} \textsuperscript{(18)} ખીલે છે.

%%K t08-c1-07-1 p
7. --- \VUgras{ધોબણ} \textsuperscript{(19)} \VUgras{નાના ઝરા}
\textsuperscript{(20)} પાસે નહાય છે, જે બે ખડકોની વચ્ચેથી ફૂટે છે, અને
\VUgras{કાંટાળાં ઝાંખરાં} \textsuperscript{(21)} તથા \VUgras{હૉથોર્નનાં
ઝાડવાં} \textsuperscript{(22)} નીચે સંતાય છે.

%%K t08-tit-3 sub
\VUsaut{5.0mm}
\VUpk{10.2pt}{40}{લણણી.}
\VUsaut{3.4mm}

%%K t08-c1-08-1 p
8. --- ગામડાનાં છોકરાં માટે ઘઉંની લણણી ઘણો મજાનો સમય છે. તેઓ આનંદથી
\VUgras{ગુલાંટો} \textsuperscript{(24)} મારે છે, જ્યાં \VUgras{પરાળના
ઢગલા} \textsuperscript{(23)} પડ્યા હોય ત્યાં ; તેઓ \VUgras{લણનારાઓને}
\textsuperscript{(26)} મદદ કરે છે, અને જ્યારે એ લોકો \VUgras{ઘઉંનું ખેતર}
\textsuperscript{(27)} પોતાનાં \VUgras{લાંબા હાથાનાં દાતરડાંથી}
\textsuperscript{(25)} વાઢે છે, જે \VUgras{ધાર કાઢવાના પથ્થર}
\textsuperscript{(30)} પર બરાબર ધારદાર કરેલાં છે, ત્યારે છોકરાં ઘઉં ભેગા
કરે છે અને તેના \VUgras{પૂળા} \textsuperscript{(31)} કે \VUgras{ભારીઓ}
\textsuperscript{(32)} બાંધે છે.

%%K t08-c1-09-1 p
9. --- ક્યારેક તેમનામાંના સૌથી મોટાને \VUgras{દાતરડાથી}
\textsuperscript{(12)} સોનેરી \VUgras{સાંઠા} \textsuperscript{(28)} કાપવા
દેવાય છે, જે દાણાથી ભરેલાં ભારે \VUgras{કણસલાં} \textsuperscript{(29)}
ધરાવે છે ; અને નાનાં છોકરાં, \VUgras{ઢોરની} \textsuperscript{(33)} ચોકી
કરતાં કરતાં, જે \VUgras{ગોચરમાં} \textsuperscript{(34)} ચરે છે અને જેની
ઉપર મોટું \VUgras{લિન્ડનનું ઝાડ} \textsuperscript{(35)} છાયા કરે છે,
પોતાની \VUgras{જાળીથી} \textsuperscript{(36)} સુંદર પતંગિયાં પકડે છે.

%%K t08-c1-09-2 p
કેટલાંક તો \VUgras{ગાડા} \textsuperscript{(37)} પર ચડી જાય છે, પૂળા
ગોઠવવા માટે, જે તેમને \VUgras{પંજેટીથી} \textsuperscript{(38)} ઉપર
ચડાવાય છે.

%%K t08-c1-10-1 p
10. --- આખા \VUgras{મેદાનમાં} \textsuperscript{(39)}, જ્યાં
\VUgras{ચંડોળ} \textsuperscript{(41)} ઊડવા લાગે છે, ચહલપહલ છવાયેલી છે.
બધા લણણીમાં રોકાયેલા છે. પણ, જ્યારે ગરીબ લોકો જૂનાં ઓજારો —
\VUgras{લાંબા હાથાનાં દાતરડાં, દાતરડાં અને ઝૂડવાની લાકડીઓ}
\textsuperscript{(40)} — વાપરવાનું ચાલુ રાખે છે, ત્યારે ધનવાન જમીનદારો
પોતાના ઘઉં કે પોતાની \VUgras{રાય} \textsuperscript{(43)} \VUgras{લણવાના
યંત્રથી} \textsuperscript{(42)} વઢાવે છે. પછી, પૂળા ખળામાં લઈ જવાની જરૂર
વિના, જગ્યા પર જ મોટી \VUgras{ગંજીઓ} \textsuperscript{(44)} કરાય છે.

%%K t08-c1-11-1 p
11. --- ત્યારે \VUgras{ઝૂડવાનું યંત્ર} \textsuperscript{(47)} લવાય છે,
જેને \VUgras{ફરતું વરાળ-એન્જિન} \textsuperscript{(45)} \VUgras{પટ્ટા}
\textsuperscript{(46)} વડે ચલાવે છે, અને એમ દાણો, જે ખેતરમાં તેને વાવ્યો
હતો તે છોડ્યા વિના જ, પરાળથી છૂટો પડે છે.

%%K t08-tit-4 sub
\VUsaut{5.0mm}
\VUpk{10.2pt}{40}{વાવાઝોડું.}
\VUsaut{3.4mm}

%%K t08-c1-12-1 p
12. --- લણણીના સમયે વારંવાર મોટાં \VUgras{વાવાઝોડાં}
\textsuperscript{(53)} થાય છે. પહેલાં \VUgras{ગર્જનાનો} ગડગડાટ દૂરથી
સંભળાય છે ; \VUgras{પશ્ચિમનો પવન} \textsuperscript{(54)} જોરથી ફૂંકાવા
લાગે છે અને હાંફતો હાંફતો \VUgras{વનરાઈનાં} \textsuperscript{(49)} સૌથી
જાડાં ઝાડ વાળી નાખે છે. કાળાં \VUgras{વાદળાં} \textsuperscript{(56)}
આકાશ પર ચડી આવે છે ; \VUgras{વીજળીના} \textsuperscript{(55)} આંજી નાખતા
વાંકાચૂકા લિસોટા તેમને ચીરે છે. \VUgras{વરસાદ} અને ક્યારેક
\VUgras{કરા} પડવા લાગે છે, જે વાઢવા માટે તૈયાર પાકને છૂંદી નાખે છે.

%%K t08-c1-13-1 p
13. --- વારંવાર વીજળી કોઈ મકાનને સળગાવી દે છે. એમ જ ગયા વર્ષે
\VUgras{પવનચક્કીની} \textsuperscript{(50)} છત રાખ થઈ ગઈ અને તેની બે
\VUgras{પાંખો} \textsuperscript{(51)} ભાંગી ગઈ.

%%K t08-c1-14-1 p
14. --- થોડા કલાક પછી શાંતિ પાછી આવે છે. ક્ષિતિજમાં ચમકતું
\VUgras{મેઘધનુષ} \textsuperscript{(52)} દેખાય છે અને કુદરત થોડી ક્ષણ માટે
તૂટેલું પોતાનું જીવન ફરી શરૂ કરે છે. \VUgras{ઝૂંપડીઓમાં}
\textsuperscript{(48)} આશરો લીધેલા ખેડૂતો ફરી કામે લાગે છે અને હિંમતથી
પોતે હમણાં જ વેઠેલું નુકસાન સરખું કરવા મથે છે.

%%K t08-tit-5 sub
\VUsaut{6.0mm}
\VUcentre{13.2pt}{90}{\textit{બીજું દૃશ્ય.}}
\VUsaut{3.6mm}

%%K t08-tit-6 sub
\VUpk{10.2pt}{40}{શિકાર. --- દ્રાક્ષની લણણી.}
\VUsaut{1.4mm}
\VUpk{10.2pt}{40}{માછીમારી.}
\VUsaut{4.0mm}

%%K t08-c2-01-1 p
1. --- પ્રદેશ પ્રમાણે, \VUgras{ઑગસ્ટના} અંતમાં કે \VUgras{સપ્ટેમ્બરની}
મધ્યમાં શિકારની મોસમ ખૂલે છે. સૂરજનાં પહેલાં કિરણોએ \VUgras{ગરોળીઓને}
\textsuperscript{(2)} તેમના દરમાંથી માંડ બહાર કાઢી હોય ત્યાં તો
\VUgras{શિકારી} \textsuperscript{(5)} પોતાના \VUgras{કૂતરાઓ}
\textsuperscript{(4)} સાથે નીકળી પડે છે.

%%K t08-c2-02-1 p
2. --- તેણે જાડા કાપડનો પોતાનો મજબૂત \VUgras{શિકારી પોશાક}
\textsuperscript{(9)} પહેર્યો છે, પગે ચામડાની \VUgras{પટ્ટીઓ} બાંધી છે,
જેનાથી તે ભીના ઘાસમાં ચાલી શકશે જ્યાં \VUgras{દેડકા}
\textsuperscript{(1)} સંતાય છે ; તેણે પોતાની \VUgras{બંદૂક}
\textsuperscript{(6)} અને પોતાની \VUgras{શિકારની થેલી}
\textsuperscript{(10)} લીધી છે, અને લો, તે નીકળી પડ્યો.

%%K t08-c2-03-1 p
3. --- પીછો કરવાના જોશમાં તે એક દ્રાક્ષવાડીની વચ્ચે આવી પહોંચે છે, જ્યાં
દ્રાક્ષની લણણી પૂરજોશમાં ચાલે છે. દ્રાક્ષ ઉગાડનારનાં છોકરાં, જે રોજ રસ્તા
પર બકરીઓ ચરાવે છે, આજે તેમને મન ફાવે તેમ ચરવા દે છે, અને ઘરડા
\VUgras{બકરાને} \textsuperscript{(3)} ખીલે બાંધ્યા પછી — જે છૂટો હોય તો
પોતાની આઝાદીનો લાભ લઈને ભાગી જ જાય — તેઓ પણ પોતાનો મજાનો ભાગ લે છે. એક
જણ \VUgras{લણણીના ટોપલામાંથી} \textsuperscript{(12)} દ્રાક્ષનો ઝૂમખો લે
છે અને \VUgras{રસ} \textsuperscript{(14)} \VUgras{પ્યાલામાં}
\textsuperscript{(13)} ટપકાવીને મજા કરે છે, જેથી મોસ્ટ (નહીં આથેલો
દ્રાક્ષરસ) બને. બીજો પોતે હમણાં જ ગૂંથેલો \VUgras{પાંદડાંનો મુગટ}
\textsuperscript{(15)} માથે પહેરે છે.

તેમની પાસે જ, નિર્લજ્જ \VUgras{ચકલીઓ} \textsuperscript{(16)} છેક દ્રાક્ષ
પીલવાના યંત્ર આગળ સુધી આવીને દ્રાક્ષના દાણા ચણી જાય છે.

%%K t08-c2-04-1 p
4. --- છોકરાં મજા કરે છે ત્યારે માણસો કામ કરે છે. \VUgras{દ્રાક્ષ
લણનારા} \textsuperscript{(18)} \VUgras{પીઠ પરના ટોપલામાં}
\textsuperscript{(17)} દ્રાક્ષ ભોંયરામાં લાવે છે ; \VUgras{પીપ
બનાવનારો} \textsuperscript{(19)} \VUgras{પીપ} \textsuperscript{(20)}
સમારે છે, \VUgras{પાટિયાં} \textsuperscript{(22)} અને
\VUgras{તળિયાં} \textsuperscript{(21)} બરાબર છે કે નહીં તે તપાસે છે, અને
ભાંગેલાં \VUgras{કડાં} \textsuperscript{(23)} બદલી નાખે છે. તેણે જ મોટાં
\VUgras{કૂંડાં} \textsuperscript{(25)} બનાવ્યાં છે, જેમાં
\VUgras{દ્રાક્ષ} \textsuperscript{(26)} આથો લાવવા માટે ઠાલવાય છે.

%%K t08-c2-05-1 p
5. --- આથો પૂરો થાય ત્યારે દારૂ ઉતારી લેવાય છે અને બાકી રહેલો ગર
\VUgras{પીલવાના યંત્ર} \textsuperscript{(28)} પર મુકાય છે, અને
\VUgras{મોટા સ્ક્રૂને} \textsuperscript{(29)} ધીમે ધીમે કસતાં જઈને રસ
નિચોવાય છે, જે \VUgras{નાના કૂંડામાં} \textsuperscript{(32)} વહે છે,
અને એમાંથી પણ \VUgras{દારૂ} \textsuperscript{(31)} મળે છે.

%%K t08-tit-8 sub
\VUsaut{6.0mm}
\VUpk{10.2pt}{40}{બિયર અને સફરજનનો દારૂ.}
\VUsaut{3.4mm}

%%K t08-c2-06-1 p
6. --- \VUgras{ભોંયરાની} \textsuperscript{(27)} દીવાલે \VUgras{હૉપની}
\textsuperscript{(30)} ડાળી ચડે છે. અહીં તે માત્ર શોભાનો છોડ છે, પણ
કેટલાક દેશોમાં, જ્યાં દ્રાક્ષવેલો બહુ ઓછો છે, ત્યાં \VUgras{બિયર}
બનાવવા માટે હૉપ ઉગાડાય છે.

%%K t08-c2-07-1 p
7. --- નાનું \VUgras{સફરજનનું ઝાડ} \textsuperscript{(33)} સફરજનથી
લદાયેલું છે, જે પાકશે ત્યારે ઉત્તમ \VUgras{સફરજનનો દારૂ} બનાવશે. પોતાના
\VUgras{પિંજરામાં} \textsuperscript{(34)} \VUgras{કૅનેરી}
\textsuperscript{(35)} અને \VUgras{ગોલ્ડફિન્ચ} \textsuperscript{(36)}
ઈર્ષાભરી નજરે એ ચકલીઓને જુએ છે, જે છૂટથી ઊડાઊડ કરે છે.

%%K t08-c2-08-1 p
8. --- નજર પહોંચે ત્યાં સુધી, ટેકરીઓના ઢોળાવ પર \VUgras{ટેકાના થાંભલા}
\textsuperscript{(40)} હારબંધ ઊભા છે, જે ભારે \VUgras{ઝૂમખાંથી}
લદાયેલા સુંદર \VUgras{દ્રાક્ષવેલાને} \textsuperscript{(39)} ટેકો આપે છે.

%%K t08-c2-08-2 p
આખું વર્ષ \VUgras{દ્રાક્ષ ઉગાડનારે} \textsuperscript{(42)} પોતાની
\VUgras{દ્રાક્ષવાડીની} \textsuperscript{(38)} સંભાળ લેવા મહેનત કરી છે,
અને હવે તેને પોતાની મહેનતનો બદલો સુંદર લણણીથી મળે છે. \VUgras{દ્રાક્ષ
લણનારીઓ} \textsuperscript{(41)} \VUgras{ખચ્ચરોથી}
\textsuperscript{(43)} ખેંચાતા ગાડા પર મૂકેલાં કૂંડાં ભરે છે, અને બધાં
આનંદમાં છે.

%%K t08-tit-9 sub
\VUsaut{5.0mm}
\VUpk{10.2pt}{40}{માછીમારી.}
\VUsaut{3.4mm}

%%K t08-c2-09-1 p
9. --- એવું જ \VUgras{ગલથી માછલી પકડનારાઓનું} \textsuperscript{(44)}
છે, જે સવારથી પાણીના કાંઠે પોતાની \VUgras{માછલી પકડવાની સોટીઓ}
\textsuperscript{(45)} લઈને બેસી ગયા છે.

%%K t08-c2-09-2 p
\VUgras{માછલી} \textsuperscript{(47)} \VUgras{ગલમાં} ફસાય છે ત્યાં જ,
જ્યાં તેઓ પોતાની \VUgras{દોરી} \textsuperscript{(46)} પાણીમાં નાખે છે,
અને \VUgras{ચારો} ફરી લગાડવાનો સમય માંડ મળે છે ત્યાં તો નવો શિકાર
દોરીના છેડે તરફડવા લાગે છે.

%%K t08-c2-09-3 p
તેમ છતાં \VUgras{જાળથી માછલી પકડનારો} \textsuperscript{(48)}, જે પોતાની
\VUgras{હોડીમાંથી} \textsuperscript{(49)} \VUgras{ફેંક-જાળ}
\VUgras{નદીની} \textsuperscript{(50)} વચ્ચે નાખે છે, વધારે માછલી પકડે છે.

%%K t08-c2-10-1 p
10. --- પાનખર સારી લટારોની ઋતુ છે : પગે ચાલીને, \VUgras{ઘોડે}
\textsuperscript{(52)}, \VUgras{સાઇકલે} \textsuperscript{(53)},
\VUgras{મોટરે} \textsuperscript{(54)}. \VUgras{રસ્તા}
\textsuperscript{(51)} સાફ છે અને છેલ્લા સારા દિવસોનો લાભ લેવાય છે, કેમ
કે લો, \VUgras{અબાબીલ} \textsuperscript{(56)} છાપરાં પર ભેગાં થવા
લાગ્યાં છે, પૂરી પાંખે વધારે ગરમ દેશો તરફ ઊડી જવા માટે. ઘરડા
\VUgras{અખરોટના ઝાડનાં} \textsuperscript{(57)} પાંદડાં પીળાં પડે છે,
\VUgras{અખરોટ} ખરે છે, અને \VUgras{ગુચ્છની} \textsuperscript{(58)} જેમ
ટોળે ગોઠવાયેલાં ઊંચાં \VUgras{ઝાડ} \VUgras{ઊંચાણ}
\textsuperscript{(59)} પર જલદી પાંદડાં વગરનાં થઈ જશે.

%%K t08-c2-11-1 p
11. --- \VUgras{સૂરજ} \textsuperscript{(60)} મોડો ઊગે છે, દિવસ ટૂંકા
થાય છે, \VUgras{સવારે} ઠીક ઠીક તીખી ઠંડી લાગવા માંડે છે, અને લો,
\VUgras{વુડકૉકનાં} \textsuperscript{(61)} ટોળાં આપણા મહેમાનગતિ વગરના
પ્રદેશો છોડી જવા લાગ્યાં છે.

\VUsaut{6.0mm}
\VUfilet{22mm}
